%% Internal planning note (Jan Streffing, co-authors), 2026-11-17: what Paper 1 does not cover, follow-up papers,
%% and open questions upstream (origins) and downstream (impacts). Not checked against new analysis; the scope list
%% states what was not done, from the paper text and docs/methods_notes.md.
%% Revised 2026-11-17 after Jan: no other seasons (Dunkelflauten are a winter phenomenon), fleet and technology out of
%% scope, cold co-occurrence and timing promoted, literature gaps of the mechanism added. References checked on Crossref.
%% Build: module load texlive/live2025-gcc-13.3.0; latexmk -pdf followup_topics.tex
\documentclass[10pt,a4paper]{article}
\usepackage[margin=1.8cm,top=1.6cm,bottom=1.6cm]{geometry}
\usepackage[T1]{fontenc}
\usepackage{lmodern}
\usepackage{enumitem}
\usepackage{titlesec}
\usepackage[hidelinks]{hyperref}
\setlist{nosep,leftmargin=1.2em}
\setlength{\parskip}{2pt}
\setlength{\parindent}{0pt}
\titlespacing*{\section}{0pt}{6pt}{2pt}
\titlespacing*{\subsection}{0pt}{4pt}{1pt}
\titleformat{\section}{\normalsize\bfseries}{\thesection}{0.6em}{}
\titleformat{\subsection}{\small\bfseries}{\thesubsection}{0.6em}{}

\begin{document}
{\large\bfseries Dunkelflauten under warming: what Paper 1 leaves out, and where to go next}\\[2pt]
{\small Internal planning note, 17 November 2026}

\section{What Paper 1 does not cover}

The paper should state these limits plainly, and none of them is analysed in it.
\begin{itemize}
\item \textbf{Timing within the winter.} Events are pooled over September to March. We did not analyse in which
      months they occur, nor whether warming shifts them towards early or late winter.
\item \textbf{Co-occurrence with cold.} We did not analyse whether events coincide with cold spells, and so with high
      demand, nor how that coincidence changes.
\item \textbf{Duration and depth.} We count spells of at least 48\,h at a matched base rate. We did not analyse the
      distribution of event lengths, the longest events, or the energy deficit during an event.
\item \textbf{Solar.} The solar change is treated as uncertain. Our runs lack aerosol-cloud interaction, and CMIP6
      disagrees on its sign.
\item \textbf{Pathway.} AWI-CM3 follows SSP5-8.5, IFS-FESOM SSP2-4.5, and CMIP6 SSP5-8.5 only. We express changes per
      kelvin and test linearity only with two CMIP6 future periods. Whether the pathway matters is not analysed.
\item \textbf{Sampling and land use.} One member per configuration, 17 and 10 winters per period at 9 and 31\,km, and
      fixed land cover in our runs.
\item \textbf{Mechanism.} We infer from the ratio of strength to slope that the day-to-day pressure deviations
      broaden, by a median of 3--5\,\%, while their strength falls less. We did not measure the length scale
      directly, did not find the cause beyond excluding the mean state and the regime pattern, and did not locate
      where around the regime pattern the gradients weaken most.
\item \textbf{Observed past.} We did not test whether the signal is already detectable in ERA5.
\end{itemize}

\section{Candidate follow-up papers}

\subsection{Cold Dunkelflauten and their timing in the winter (first priority)}
The most expensive events combine little wind and sun with cold and high demand \cite{vanderWiel2019}. Blocking
favours both, since it brings calm air and, from the east, cold air. Questions: how often do low-wind events
coincide with cold spells today, and in which regime; does warming change that joint probability, given that cold
extremes in northern mid-latitudes warm faster than the mean \cite{Screen2014}; and do events shift within the winter,
for example towards December and January, when solar output is lowest and demand highest. The same events answer the
duration and depth questions of Sect.~1. Events and regimes exist for all runs, while daily 2\,m temperature still
has to be extracted (the repository holds only box means for some runs and monthly values for CMIP6). This is the
cheapest follow-up and connects directly to demand, without entering fleet or system modelling.

\subsection{Does the pathway matter? (climate journal)}
Test whether the change per kelvin depends on global warming alone. CMIP6 provides several scenarios with daily wind
and pressure, idealised 1\,\%\,CO$_2$ and abrupt-4$\times$CO$_2$ runs, and stabilisation runs after zero emissions.
Candidate pathway effects are the regional aerosol decline, the rate of warming against ocean adjustment (AMOC,
North Atlantic warming hole), the timing of Arctic sea-ice loss, and overshoot. The CMIP7 runs of AWI-ESM3 with
several pathways would add the same model family and, ideally, aerosol-cloud interaction for the solar side.

\subsection{A constraint on the projections, if a first check supports it}
That weather systems broaden under warming is known. Models robustly enlarge the eddy length scale
\cite{Kidston2010}, linked to the poleward shift of the jets \cite{Kidston2011}. Storm tracks feel opposing
influences from a weaker low-level and a stronger upper-level temperature gradient \cite{Shaw2016}, Arctic amplification
damps temperature variance \cite{Screen2014}, and blocking projections over Europe remain uncertain \cite{Woollings2018}.
Why the length scale grows is unsettled, but answering that needs idealised experiments and theory beyond our
strength, so a mechanism paper of our own is not recommended. Paper 1 already contains the new part, the link to the
low-wind tail. What could still carry a paper is a constraint. If the CMIP6 models that broaden more also gain more
low-wind events, the spread in broadening explains part of the spread in the projections, and ERA5 since 1950 might
show how much broadening is realistic. The first check is cheap, since broadening and event change per model exist
for all 20 models. Without a correlation, this becomes a question in Sect.~3.

\section{Further questions}

\textbf{Upstream, towards the origins}
\begin{itemize}
\item Why does the length scale grow? Candidate explanations are a larger deformation radius from static stability
      and tropopause height, weaker low-level baroclinicity, and latent heating, and none is settled.
\item Is the broadening a dry dynamical response or a moist one? Aquaplanet or uniform-SST experiments with OpenIFS
      could separate them.
\item Where around the regime pattern do the gradients weaken most, and why do the hemispheres differ at equal mean
      gradient?
\item How much comes from Arctic amplification and sea-ice loss? Prescribed sea-ice experiments (as in PAMIP) would test it.
\item Why does EC-Earth3, from the same IFS family, weaken $G$ more than any other CMIP6 model, while our IFS-based
      runs spread across the CMIP6 range? A shared code base with different results is a controlled test.
\item Do sudden stratospheric warmings, which favour blocking and cold spells, change in a way that interacts with
      the broadening?
\item When does the signal emerge from variability?
\end{itemize}

\textbf{Downstream, towards the impacts}
\begin{itemize}
\item Does the drop in heating demand from warmer winters offset the rise in events, or does electrified heating
      keep cold Dunkelflauten critical?
\item Does the broadening make low-wind spells more coherent across Europe, and so harder to balance across borders?
\item Are blocking-related Dunkelflauten predictable at subseasonal range, and does warming change that?
\item Do the same weaker gradients also reduce damaging wind storms?
\item What do the results imply for storage needs? This is a question for energy-system partners, with our series as
      input.
\end{itemize}

{\small
\begin{thebibliography}{9}\setlength{\itemsep}{0pt}
\bibitem{Kidston2010} Kidston, J., Dean, S. M., Renwick, J. A., and Vallis, G. K. (2010), Geophys. Res. Lett. 37, doi:10.1029/2009GL041615.
\bibitem{Kidston2011} Kidston, J., Vallis, G. K., Dean, S. M., and Renwick, J. A. (2011), J. Climate 24, 3764, doi:10.1175/2010JCLI3738.1.
\bibitem{Screen2014} Screen, J. A. (2014), Nat. Clim. Change 4, 577, doi:10.1038/nclimate2268.
\bibitem{Shaw2016} Shaw, T. A., et al. (2016), Nat. Geosci. 9, 656, doi:10.1038/ngeo2783.
\bibitem{vanderWiel2019} van der Wiel, K., et al. (2019), Renew. Sustain. Energy Rev. 111, 261, doi:10.1016/j.rser.2019.04.065.
\bibitem{Woollings2018} Woollings, T., et al. (2018), Curr. Clim. Change Rep. 4, 287, doi:10.1007/s40641-018-0108-z.
\end{thebibliography}}

\end{document}
